\documentclass[conference]{IEEEtran}
\IEEEoverridecommandlockouts

\usepackage{cite}
\usepackage{amsmath,amssymb,amsfonts}
\usepackage{algorithmic}
\usepackage{graphicx}
\usepackage{textcomp}
\usepackage{xcolor}
\usepackage{pifont}
\usepackage{booktabs}
\usepackage{tabularx}
\usepackage{array}

\newcommand{\cmark}{\ding{51}} 
\newcommand{\xmark}{\ding{55}} 

\def\BibTeX{{\rm B\kern-.05em{\sc i\kern-.025em b}\kern-.08em
    T\kern-.1667em\lower.7ex\hbox{E}\kern-.125emX}}
\begin{document}

\title{ORBIT: Objective Reasoning Benchmark for ITOps Automation
}

% \author{\IEEEauthorblockN{1\textsuperscript{st} Given Name Surname}
% \IEEEauthorblockA{\textit{dept. name of organization (of Aff.)} \\
% \textit{name of organization (of Aff.)}\\
% City, Country \\
% email address or ORCID}
% \and
% \IEEEauthorblockN{2\textsuperscript{nd} Given Name Surname}
% \IEEEauthorblockA{\textit{dept. name of organization (of Aff.)} \\
% \textit{name of organization (of Aff.)}\\
% City, Country \\
% email address or ORCID}
% \and
% \IEEEauthorblockN{3\textsuperscript{rd} Given Name Surname}
% \IEEEauthorblockA{\textit{dept. name of organization (of Aff.)} \\
% \textit{name of organization (of Aff.)}\\
% City, Country \\
% email address or ORCID}
% \and
% \IEEEauthorblockN{4\textsuperscript{th} Given Name Surname}
% \IEEEauthorblockA{\textit{dept. name of organization (of Aff.)} \\
% \textit{name of organization (of Aff.)}\\
% City, Country \\
% email address or ORCID}
% \and
% \IEEEauthorblockN{5\textsuperscript{th} Given Name Surname}
% \IEEEauthorblockA{\textit{dept. name of organization (of Aff.)} \\
% \textit{name of organization (of Aff.)}\\
% City, Country \\
% email address or ORCID}
% \and
% \IEEEauthorblockN{6\textsuperscript{th} Given Name Surname}
% \IEEEauthorblockA{\textit{dept. name of organization (of Aff.)} \\
% \textit{name of organization (of Aff.)}\\
% City, Country \\
% email address or ORCID}
% }

\maketitle

\begin{abstract}
This document is a model and instructions for \LaTeX.
This and the IEEEtran.cls file define the components of your paper [title, text, heads, etc.]. *CRITICAL: Do Not Use Symbols, Special Characters, Footnotes, 
or Math in Paper Title or Abstract.
\end{abstract}

\begin{IEEEkeywords}
component, formatting, style, styling, insert
\end{IEEEkeywords}

\section{Introduction}\label{sec1}

% current landscape of the domain of the study
% Problem statement & Motivation
% How is our study a solution to the problem statement identified? Implications and Applications of the Study

The rapid proliferation of complex, distributed IT infrastructures has elevated Application Performance Management (APM) and Site Reliability Engineering (SRE) to mission-critical disciplines, in which operational teams must continuously reason across heterogeneous telemetry streams, system logs, and dynamic configuration states under stringent time constraints. The convergence of AIOps, defined by the systematic integration of machine learning into ITOps workflows, with large language models (LLMs) has catalyzed a fundamental paradigm shift toward intelligent, automated operational reasoning. Empirical evidence indicates that 60\% of enterprises are projected to adopt DevOps analytics tooling, reflecting broad institutional recognition that manual incident triage and workflow bookkeeping are computationally insufficient at scale. Currently, frontier models have demonstrated consistent performance across many tasks, including document summarization, reviews, coding, debugging, test case generation, issue resolution, and general reasoning. 

Despite these advances, a structurally persistent gap persists between offline benchmark performance and real-world operational deployment reliability. Evaluations, including SWE-bench, ITBench, and OpsEval, are respectively code-centric, execution-centric, or constrained to single-turn question-answering paradigms, none of which capture the diagnostic, multi-step inferential structure that is intrinsic to benchmarking the reasoning aspects required to execute ITOps workflows. Critically, no existing benchmark simultaneously provides multi-step reasoning annotations grounded in realistic operational contexts, while pervasive privacy restrictions further constrain real-world data accessibility, rendering targeted LLM capability assessment fundamentally infeasible.

This paper directly addresses the identified gap by introducing a reasoning benchmark explicitly constructed for IT automation, systematically grounded in operational wikis and practitioner forums, sources that authentically encode the diagnostic heuristics and inferential patterns employed by domain experts. By synthesizing knowledge from Application Monitoring Management tools (APMs) such as Instana, Datadog, Dynatrace, etc., documentation, and real-world incident-resolution discussions on Reddit and Stack Overflow, the proposed dataset provides multi-step reasoning annotations anchored in executable operational contexts, enabling rigorous LLM evaluation under realistic conditions. This resource facilitates targeted diagnosis of model failure modes, particularly long-context localization and multi-hop causal inference, advancing self-aware, auto-healing system development, while simultaneously serving as a structured fine-tuning foundation for production-oriented, operationally grounded AI reasoning systems.

% Taxonomy diagram of the dataset

% Contributions

This paper makes the following contributions:
\begin{itemize}

    \item A well-curated reasoning benchmark dataset for ITOps, systematically grounded in operational wikis and practitioner forums, providing multi-step reasoning annotations anchored in authentic diagnostic and incident-resolution contexts.
    
    \item An extensive benchmarking study evaluating large language models across operationally grounded multi-step reasoning tasks, enabling targeted diagnosis of model failure modes, particularly long-context localization and multi-hop causal inference, under realistic IT automation conditions.
    
\end{itemize}

% Organization of the research study

The paper is organized as follows. Section~\ref{sec2} reviews related works; Section~\ref{sec3} details the dataset construction methodology; Section~\ref{sec4} details the motivation of the study and throws light on the preliminaries required for the study; Section~\ref{sec5} describes the annotation system implementation; Section~\ref{sec6} presents the experimental setup and benchmarking results; and Section~\ref{sec7} concludes with key findings and future directions.

\section{Related Works}\label{sec2}

% The Evolution of ITOps to AIOps
% The Role of Generative AI in Automation
% Evaluating AI Reasoning in Operations
% Dataset Construction Methodologies
% Research Gap

The transition from traditional ITOps to AIOps has introduced intelligent, automated paradigms that fundamentally reshape infrastructure management. \cite{dang2019aiops} synthesize real-world AIOps experiences from Microsoft, noting that 60\% of firms will adopt DevOps analytics by 2024, while identifying persistent limitations including poor data quality, missing ground-truth labels, and complex system dependencies. This foundational work advocates for continuous academia-industry collaboration to build self-aware, auto-healing systems. Complementing this, \cite{zheng2023survey} systematically surveyed 134 papers on Code LLMs, demonstrating that code-specialized models consistently outperform general-purpose counterparts, establishing the technical foundation for deploying intelligent operational tooling at scale.

Generative AI has demonstrated strong potential for operational automation, yet a recurring gap between offline benchmark performance and real-world operational reliability persists across approaches. \cite{sacco2025intent} showed this through Kubernetes manifest generation, where LoRA accelerated training by 3× and reduced memory by 45\%, while full fine-tuning retained superior task accuracy, revealing that efficiency optimizations trade against fidelity. \cite{saha2024sequential} advanced sequential API evaluation through GraphQLRestBench, yet non-executable settings limit operational transferability. \cite{shan2025rca} introduced RCACopilot, where few-shot retrieval achieved 0.95 BERTScore F1 against 0.66 zero-shot, confirming that grounding consistently improves reliability but remains architecturally contingent on retrieval quality.

Current benchmarks consistently expose deep, structurally specific failures in AI reasoning under realistic operational conditions. \cite{jimenez2023swe} introduced SWE-bench across 2,294 real GitHub issues, where Claude 2 resolved only 1.96\%, reflecting empirically observed failures in long-context code localization. \cite{jha2025itbench} deployed ITBench across 94 Kubernetes-grounded tasks, where GPT-4o resolved 0\% of FinOps and only 13.8\% of SRE scenarios. Execution non-determinism from live telemetry and incomplete observability were identified as the primary failure drivers. These results empirically demonstrate that operational reasoning demands multi-step, execution-aware judgment that current models struggle to sustain across realistic task conditions.

These empirical failures are linked to a primary bottleneck at the dataset level. Robust dataset construction is imperative for reliably assessing AI. LogPPT \cite{le2023log} achieves 0.923 Group Accuracy across sixteen log datasets through an adaptive random sampling algorithm for few-shot learning, and GraphQLRestBench \cite{saha2024sequential} constructs sequential API pipelines, yet both remain narrowly scoped to single-step operational sub-tasks. \cite{liu2025opseval} further confirms that privacy restrictions severely constrain real-world data access. Critically, no existing dataset provides multi-step reasoning annotations grounded in authentic operational contexts. Because benchmark failures trace specifically to execution-aware, multi-step reasoning, and current datasets do not capture that reasoning structure, the design of annotations represents a primary bottleneck to further progress.

\begin{table*}[t]
\centering
\caption{Comparison of existing benchmarks and ORBIT across key operational reasoning requirements.}
\label{tab:benchmark_comparison}
\small
\begin{tabular}{lccccc}
\toprule
\textbf{Benchmark} &
\textbf{Operational} &
\textbf{Multi-Step} &
\textbf{Execution-Aware} &
\textbf{Evidence} &
\textbf{Reasoning} \\
&
\textbf{Grounding} &
\textbf{Inference} &
\textbf{Reasoning} &
\textbf{Traceability} &
\textbf{Annotations} \\
\midrule
SWE-bench~\cite{jimenez2023swe}      & \xmark & \cmark & \xmark & \cmark & \xmark \\
ITBench~\cite{jha2025itbench}        & \cmark & \cmark & \cmark & \xmark & \xmark \\
OpsEval~\cite{liu2025opseval}        & \cmark & \xmark & \cmark & \xmark & \xmark \\
LogPPT~\cite{le2023log}              & \cmark & \xmark & \xmark & \xmark & \xmark \\
GraphQLRestBench~\cite{saha2024sequential} & \xmark & \cmark & \xmark & \cmark & \xmark \\
\midrule
\textbf{ORBIT (Proposed Dataset)}                & \textbf{\cmark} & \textbf{\cmark} & \textbf{\cmark} & \textbf{\cmark} & \textbf{\cmark} \\
\bottomrule
\end{tabular}
\end{table*}

The observations summarized in Table~\ref{tab:benchmark_comparison} motivate four design requirements for an operational reasoning benchmark:

\begin{enumerate}
    \item \textbf{Operational Grounding:} Examples must be derived from authentic infrastructure management, observability, cloud operations, and reliability engineering workflows.

    \item \textbf{Multi-Hop Reasoning:} Tasks should require models to synthesize information across multiple evidence sources rather than retrieve isolated facts.
    
    \item \textbf{Evidence Traceability:} Every reasoning instance should maintain explicit provenance linking conclusions to supporting operational evidence.
    
    \item \textbf{Technology Diversity:} The benchmark should cover a broad spectrum of operational technologies, including cloud platforms, infrastructure-as-code frameworks, observability tools, orchestration systems, and reliability engineering practices.
\end{enumerate}


\textbf{Research Gap:} Despite notable advances, no existing resource simultaneously combines operational knowledge, multi-step reasoning structure, and realistic execution context. ITBench \cite{jha2025itbench}, OpsEval \cite{liu2025opseval}, and SWE-bench \cite{jimenez2023swe} are respectively execution-centric, QA-style, and code-centric. Adapting these resources is insufficient because their annotation schemas do not capture the diagnostic, multi-step inference characteristic of real ITOps workflows. Furthermore, \cite{dang2019aiops} and \cite{sacco2025intent} highlight compounding limitations, including missing ground-truth labels, scarce datasets, and ongoing reasoning collapse in models. This paper directly addresses this gap by constructing a reasoning dataset grounded in operational wikis and practitioner forums, enabling targeted evaluation of LLM analytical capabilities within operationally grounded contexts.

\section{Motivation \& Preliminaries}\label{sec3}

Modern Information Technology Operations (ITOps) environments have evolved into deeply interconnected, cloud-native ecosystems composed of microservices, orchestration frameworks, distributed storage systems, observability pipelines, and dynamically scaling infrastructure. Consequently, operational failures rarely arise from isolated components; instead, they propagate across heterogeneous service dependencies, requiring operators to reason simultaneously about telemetry, configuration states, deployment histories, and execution traces. For example, diagnosing latency escalation in Kubernetes may involve correlating pod restarts, network policies, service mesh bottlenecks, and infrastructure resource exhaustion under strict temporal constraints. Existing evaluation benchmarks primarily assess isolated coding or question-answering capabilities, which are insufficient to represent realistic operational reasoning requirements.

In contemporary enterprise infrastructure management, ITOps refers to the administration, monitoring, optimization, and maintenance of distributed computing systems that support organizational workloads and digital services. Application Performance Management (APM) platforms, including Instana, Datadog, and Dynatrace, continuously collect telemetry streams comprising logs, traces, metrics, events, and dependency graphs to enable observability-driven analytics and operational diagnostics. Site Reliability Engineering (SRE) extends these principles through reliability-focused automation, incident response coordination, service-level objective enforcement, and resilience engineering. AIOps further integrates machine learning and large language models into operational workflows for anomaly detection, root-cause analysis, remediation assistance, and automated infrastructure reasoning across heterogeneous operational artifacts and telemetry modalities.

The operational complexity of modern infrastructure introduces reasoning challenges that fundamentally differ from those in conventional software engineering benchmarks. Real-world incidents frequently require multi-hop inferential reasoning across heterogeneous evidence sources before reliable remediation decisions can be established. For instance, Terraform misconfigurations may indirectly trigger Kubernetes scheduling instability, subsequently causing cascading application failures, elevated latency, and distributed service degradation. Similarly, diagnosing cloud resource exhaustion often requires simultaneous reasoning over deployment timelines, autoscaling policies, infrastructure metrics, historical telemetry, and practitioner documentation. Existing benchmarks such as SWE-bench, ITBench, and OpsEval reveal persistent failures in long-context localization, sequential causal inference, and execution-aware reasoning, highlighting a substantial gap between offline benchmark performance and real-world operational reliability.

The motivation for ORBIT emerges from the persistent mismatch between benchmark success and real-world operational reliability observed in existing evaluation frameworks. Benchmarks, including SWE-bench, ITBench, and OpsEval, reveal substantial failures in long-context localization, execution-aware inference, and operational reasoning under realistic deployment conditions. For instance, infrastructure failures involving Terraform misconfigurations, Kubernetes networking anomalies, or cloud resource exhaustion often require multi-stage inferential reasoning that simultaneously combines documentation, practitioner discourse, and historical telemetry. However, privacy restrictions and limited access to production operational data significantly constrain the construction of datasets for realistic evaluation. ORBIT addresses this gap by introducing a benchmark grounded in operational wikis, practitioner forums, and executable infrastructure reasoning contexts, specifically designed to evaluate multi-step analytical reasoning in IT automation environments.

\section{\sc{ORBIT Construction}}\label{sec4}

Constructing ORBIT, a high-quality ITOps reasoning dataset, requires a rigorous, systematic approach to data collection, processing, and validation. This section delineates the two primary phases of our dataset curation pipeline: (1) Dataset Generation, which encompasses the automated and semi-automated creation of reasoning instances from heterogeneous knowledge sources, and (2) Multi-Annotation, which establishes inter-annotator agreement and ensures label quality through iterative consensus mechanisms. Together, these phases ensure that the resulting dataset exhibits both breadth of coverage across ITOps domains and depth of annotated reasoning traces suitable for training and evaluating large language models.

\subsection{Data Generation \& Collection}

The data generation workflow was implemented using a dedicated annotation platform, \textit{Dataset Annotator}, engineered to construct operational reasoning benchmarks for ITOps automation. The platform employed Next.js 16 with React 19 and adopted a dual backend architecture integrating Neon PostgreSQL with \texttt{pgvector} for scalable semantic retrieval and SQLite with \texttt{sqlite-vec} for localized experimentation. Gemini API services supported semantic embedding generation and AI-assisted draft synthesis, enabling unified dataset authoring, provenance tracking, semantic retrieval, and annotation management. A version-controlled relational schema preserved complete reasoning lifecycles by storing structured question-answer content, reasoning traces, conclusions, and coverage metadata in the \texttt{dataset\_examples} and \texttt{dataset\_states} tables, which are immutable.

To establish a semantically grounded operational knowledge repository, the system incorporated a dedicated ingestion pipeline accessible through the \texttt{/api/itops/ingest} interface. Retrieved HTML documents originating from operational wikis, practitioner resources, and APM documentation were normalized using a Docling-style parser that removed non-semantic interface components while preserving diagnostically relevant structures, including headings, code blocks, lists, and tables. Extracted content was segmented using sentence-aware chunking with configurable token overlap to maintain contextual continuity across adjacent segments. Subsequently, each chunk was encoded into a 3072-dimensional embedding using the Gemini embedding model and indexed in the vector database for retrieval-augmented generation during synthesis.

AI-assisted dataset synthesis was orchestrated through the \texttt{/api/generate} endpoint, which accepted structured prompts containing user objectives, task specifications, retrieved grounding evidence, and topic-level constraints. The platform supported multiple benchmark formats, including multiple-choice reasoning tasks, open-ended question answering, and multi-step inferential chains, while generated outputs followed a canonical JSON schema containing questions, answer candidates, reasoning explanations, thought processes, and concept coverage metadata. Operational relevance was preserved through curated technology registries spanning Kubernetes, Terraform, AWS, Azure, Instana, and Turbonomic, along with stop-topic filtering rules. Finalized datasets were exported via normalized JSON pipelines that support provenance preservation, duplicate detection, transactional consistency, benchmarking, and downstream fine-tuning.

\subsection{Multi-Annotation For the Dataset}

To ensure annotation consistency and benchmark reliability, we developed a dedicated multi-annotator evaluation framework using a Streamlit-based labeling interface integrated with PostgreSQL-backed storage. The framework supported inter-annotator agreement analysis, consensus formation, and collaborative validation of reasoning instances generated by the Dataset Annotator. Multiple independent annotators reviewed identical examples under standardized presentation conditions, enabling systematic identification of ambiguous questions, incorrect reasoning traces, and insufficiently grounded answers while reducing individual labeling bias. The underlying relational schema comprised interconnected dataset\_items, annotators, and annotations tables that store canonical benchmark content, annotator metadata, labeling outcomes, corrected answers, explanatory notes, and provenance information, using JSONB representations efficiently.

The labeling workflow employed lightweight session-based annotator authentication within the Streamlit interface and presented dataset items using deterministic randomized sequencing generated from pseudo-random seeds derived from annotator identities and dataset size. Each instance included operational questions, answer candidates, expected responses, reasoning chains, explanatory traces, and grounding references through expandable interface panels. Annotators evaluated semantic correctness, reasoning coherence, and operational validity before assigning acceptance or rejection verdicts, with corrections and subsequent review required. This standardized evaluation procedure ensured reproducibility, fairness, and consistency across distributed annotation sessions while facilitating the detection of operational inconsistencies, reasoning errors, ambiguities, and insufficiently validated benchmark examples throughout the assessment.

Annotation persistence employed PostgreSQL upserts, allowing annotators to revise existing labels without generating duplicate records while maintaining comprehensive provenance metadata and timestamped histories. Session-level statistics continuously tracked completed annotations, remaining workload, and acceptance or rejection distributions, whereas aggregation utilities computed per-item and per-annotator consistency metrics using database-level summary queries. Synchronization scripts regularly update benchmark entries from the generation pipeline into the annotation repository, ensuring alignment between evolving dataset revisions and downstream validation workflows. The framework also supported iterative consensus formation through majority voting or expert adjudication, enabling longitudinal analysis of disagreement patterns, annotation stability, reviewer calibration, and reasoning ambiguities.

\subsection{Snorkel/UQ - How hard is the dataset}

\subsection{ORBIT Dataset \& Statistics}
\label{sec:dataset_stats}

\begin{table*}[t]
\centering
\caption{Representative ORBIT dataset samples spanning diverse operational domains, reasoning skills, and grounding sources. Questions, options, and reasoning summaries are abbreviated for readability.}
\label{tab:orbit_dataset_preview}

\footnotesize
\setlength{\tabcolsep}{3pt}
\renewcommand{\arraystretch}{1.15}

\begin{tabularx}{\textwidth}{
|p{0.8cm}
|p{1.3cm}
|p{1.2cm}
|p{1.8cm}
|p{1.7cm}
|X
|X
|p{0.55cm}
|X|}
\hline

\textbf{ID} &
\textbf{Category} &
\textbf{Domain} &
\textbf{Reasoning Skill} &
\textbf{Grounding Source} &
\textbf{Question} &
\textbf{Options} &
\textbf{Ans.} &
\textbf{Reasoning Summary} \\
\hline

INST-101 &
Trouble-shooting &
Observabi-lity &
Multi-hop Causal Reasoning &
Community Discussion &
PHP-FPM requests intermittently hang for 20--60 seconds after output generation. What monitoring-related root cause was identified? &
A: UDP metric corruption; B: TCP metric transmission contention; C: OOM-killed collector; D: Missing OTel correlation &
B &
Correlates application traces, packet captures, and monitoring behavior. Resource contention during Instana TCP metric transmission delayed request completion despite PHP execution having finished. \\
\hline

DP-00029 &
Product Understanding &
Instana &
Semantic Interpretation &
Product Documentation &
Within the Instana documentation hierarchy, what functional purpose does \texttt{.instana} serve? &
A: IBM Instana observability platform; B: Deprecated parameter; C: Monitoring latency field; D: Authentication schema &
A &
Requires understanding product taxonomy and documentation structure. The entry represents the Instana platform entity rather than a schema attribute or configuration parameter. \\
\hline

DP-00041 &
Schema Extraction &
AWS &
Structured Retrieval &
API Reference &
What is the complete field structure of \texttt{AwsEc2Instance}? &
A--C: Incorrect schema definitions; D: \texttt{aws\_resource}, \texttt{monitoring}, \texttt{private\_ip}, \texttt{tenancy}, \texttt{virtual\_machine} &
D &
Requires retrieval of the exact schema definition from AWS documentation and differentiation from related infrastructure resource types. \\
\hline

DP-00047 &
Schema Extraction &
IAM &
Attribute Reasoning &
API Reference &
Which fields belong to the \texttt{AwsIamRole} object definition? &
A: Infobox template; B: Generic type definition; C: 10-field IAM role schema; D: EC2 instance type schema &
C &
Requires identifying IAM-specific attributes, including permissions boundaries, role metadata, and session configurations, while rejecting unrelated AWS schemas. \\
\hline

DP-00025 &
Documenta-tion Navigation &
GCP &
Classification &
Product Documentation &
Based on the documentation hierarchy, what type of entry is \texttt{.gcp}? &
A: Category page; B: Deprecated parameter; C: Monitoring component; D: Google Cloud Platform infrastructure resources &
D &
Requires semantic interpretation of documentation organization and recognition of GCP as a cloud infrastructure product category. \\
\hline

DP-00056 &
Configura-tion Analysis &
Security Groups &
Structured Reasoning &
API Reference &
According to AWS schema documentation, what are the fields of \texttt{AwsSecurity GroupRule}? &
A: \texttt{aws\_resource}, \texttt{security\_rule}; B: MappingStatus definition; C: GCP type definition; D: S3 bucket schema &
A &
Requires locating the correct security-group rule schema and differentiating it from related networking, storage, and metadata object definitions. \\
\hline

\end{tabularx}
\end{table*}

The final ORBIT benchmark comprises 3,350 operational reasoning instances curated via the human-in-the-loop pipeline described previously (sample of the dataset included in Table.~\ref{tab:orbit_dataset_preview}). Each instance contains a problem statement, candidate responses (when applicable), a verified answer, structured reasoning traces, and provenance links to the original grounding sources. The dataset's taxonomy is described in Fig.~\ref{}, and the benchmark is intentionally designed to evaluate operational decision-making rather than isolated factual recall, requiring models to integrate troubleshooting knowledge, infrastructure semantics, system dependencies, and procedural reasoning.

\begin{figure*}
    \centering
    \includegraphics[width=0.8\linewidth]{itops_taxonomy.png}
    \caption{Taxonomy of the ORBIT dataset. The dataset is organized along three complementary dimensions: (1) question type distribution capturing the nature of reasoning tasks, (2) source type capturing the provenance of knowledge, and (3) knowledge quality profile highlighting accuracy and solution depth characteristics across sources.}
    \label{taxonomy}
\end{figure*}

ORBIT covers three complementary question categories. Multiple-choice operational reasoning questions constitute the majority of the benchmark (75.0\%), followed by open-ended reasoning tasks (21.0\%) and highly specialized technical troubleshooting scenarios (4.0\%). This distribution reflects the composition of real-world operational workflows, where engineers frequently alternate between diagnosis, root-cause analysis, system interpretation, and decision selection under incomplete information.

Across the benchmark, questions contain an average of 14.1 tokens and approximately 3.2 domain-specific technical entities, indicating substantial operational content density. Each sample is grounded in one or more external evidence sources, with an average of 1.03 supporting references per instance. Approximately 14.1\% of examples incorporate multiple evidence sources, requiring cross-source information synthesis. The benchmark further incorporates both positive and adversarial examples, resulting in a class distribution of 76.7\% correct instances and 23.3\% adversarial instances, enabling evaluation of calibration, robustness, and uncertainty-aware reasoning.

\begin{table}[t]
\centering
\caption{Key structural and semantic characteristics of ORBIT. Values are reported as mean $\pm$ standard deviation.}
\label{tab:orbit_feature_statistics}
\begin{tabular}{lcc}
\toprule
\textbf{Metric} & \textbf{Value} & \textbf{Unit} \\
\midrule
Question length & $14.11 \pm 5.47$ & tokens \\
Question complexity score & $1.49 \pm 0.30$ & normalized score \\
Technical terms per question & $3.19 \pm 1.72$ & count \\
Average answer length & $7.87 \pm 7.06$ & tokens \\
Grounding sources per sample & $1.03 \pm 0.50$ & count \\
Multi-source instances & $14.1$ & \% \\
Semantic consistency score & $0.944 \pm 0.041$ & score [0,1] \\
Semantic entropy & $0.234 \pm 0.666$ & entropy score \\
\bottomrule
\end{tabular}
\end{table}

The Table~\ref{tab:orbit_feature_statistics} reports key structural and semantic characteristics of the dataset. Together, these statistics demonstrate that ORBIT captures the diversity, grounding requirements, and reasoning complexity encountered in practical IT operations environments.


\section{Methodology}\label{sec5}

To systematically construct ORBIT as described in Fig.~\ref{methodology}, we developed a human-in-the-loop annotation framework that integrates retrieval-augmented knowledge grounding, model-assisted example generation, expert validation, and provenance tracking within a unified curation pipeline. The framework was designed to address the unique challenges of ITOps reasoning benchmarks, where reliable evaluation requires multi-step inferential reasoning grounded in authentic operational knowledge rather than isolated question-answer pairs. Accordingly, the methodology is guided by four core principles: (i) complete provenance and version control for every benchmark instance, (ii) semantic retrieval over curated operational knowledge repositories, (iii) grounding against authoritative practitioner discourse, and (iv) controlled diversity to ensure broad coverage across operational concepts and technologies. The overall workflow transforms heterogeneous operational resources into validated reasoning examples suitable for benchmarking large language models in realistic IT automation scenarios.

\begin{figure*}
    \centering
    \includegraphics[scale = 0.4]{methodology_ITOPS.png}
    \caption{Workflow Diagram for ORBIT.}
    \label{methodology}
\end{figure*}

\subsection{Version-Controlled Annotation Framework}

Each benchmark instance is represented as an immutable example identifier associated with an append-only sequence of annotation states. Every modification, whether performed by an annotator or an AI-assisted generation component, produces a new version containing the question, reasoning trace, answer, metadata, and revision information. This design preserves the complete evolution of an example and enables transparent auditing of all annotation decisions. Furthermore, each version maintains explicit links to the operational resources used during its creation, establishing traceability between the generated reasoning process and its supporting evidence. Such provenance preservation is essential for validating the reliability and reproducibility of multi-step reasoning annotations.

\subsection{Retrieval-Grounded Operational Knowledge Base}

ORBIT is grounded in a curated repository of operational knowledge derived from APM documentation, infrastructure management guides, cloud-native technology references, and operational wikis. Source documents undergo HTML-aware normalization to preserve diagnostically relevant structures, including code snippets, configuration artifacts, procedural instructions, tables, and hierarchical document organization. The processed content is segmented into semantically coherent chunks and encoded using dense vector representations before being indexed within a retrieval system. During annotation, natural-language queries issued by annotators retrieve the most relevant passages through semantic similarity search, providing evidence candidates that serve as grounding material for benchmark construction. This retrieval-centric design ensures that examples reflect realistic operational scenarios and domain-specific reasoning patterns encountered in production environments.

\subsection{Practitioner-Grounded Evidence Integration}

While operational documentation provides structured technical knowledge, real-world incident diagnosis frequently relies on practitioner experience and community-driven troubleshooting strategies. To capture this dimension, ORBIT incorporates Stack Overflow as an external source of operational ground truth. Highly ranked practitioner responses are retrieved and associated with relevant operational topics, enabling benchmark examples to reflect expert diagnostic reasoning and remediation practices. Within the annotation pipeline, practitioner discussions serve as authoritative reference points, whereas retrieved documentation serves as supporting evidence that must be validated and synthesized rather than reproduced directly. This separation between documentation-based evidence and practitioner-derived expertise encourages the creation of reasoning chains that more accurately resemble real-world operational decision-making processes.

\subsection{Model-Assisted Benchmark Construction and Diversity Enforcement}

To improve annotation efficiency while maintaining quality, the framework supports both manual authoring and AI-assisted example generation. Given a target task specification, user objective, and selected grounding sources, a frontier large language model generates structured benchmark drafts that may take the form of multiple-choice reasoning tasks, open-ended question-answer pairs, or multi-step diagnostic workflows. To prevent excessive focus on frequently occurring technologies or operational patterns, a concept registry continuously tracks previously used concepts across the dataset. Concepts already represented in the benchmark are incorporated into generation prompts as exclusion constraints, while annotator-selected topics are enforced as mandatory inclusion constraints. This mechanism promotes broad coverage across cloud infrastructure, observability platforms, reliability engineering practices, deployment systems, and operational troubleshooting scenarios, thereby reducing topical redundancy and improving benchmark diversity.

\subsection{Human Validation and Benchmark Finalization}

Regardless of whether an example originates from manual authoring or AI-assisted generation, every benchmark instance undergoes expert review before inclusion in ORBIT. Annotators evaluate semantic correctness, reasoning coherence, operational validity, and evidence alignment, making revisions where necessary. Retrieved operational documents and practitioner references are consolidated into a unified evidence set, which is simultaneously provided to the generation process and stored in the example's provenance record. Consequently, the evidence visible during benchmark construction is identical to the evidence retained for future auditing and reproducibility. Following validation, examples are exported in a normalized JSON format that includes the question, answer, reasoning trace, grounding sources, provenance metadata, and concept coverage information, forming the final benchmark corpus used for downstream evaluation and analysis of large language models in IT automation environments.

\subsection{Experimentation}

\section{Results \& Discussion}\label{sec6}
\section{Conclusion}\label{sec7}

\bibliographystyle{IEEEtran}
\bibliography{ref}

\end{document}

